\documentclass[12pt]{article}
\usepackage{amsmath, amssymb, amsthm}
\usepackage{geometry}
\usepackage{array}
\usepackage{enumitem}
\usepackage{titlesec}
\usepackage[dvipsnames, table]{xcolor}
\usepackage{fancyhdr}
\usepackage{tcolorbox}
\usepackage{microtype}
\usepackage{multicol}
\usepackage{booktabs}
\usepackage{multirow}
\usepackage{titletoc}
\usepackage{tikz}

\tcbuselibrary{skins, breakable}

\newtcolorbox{examplebox*}[1][]{
  enhanced, breakable,
  colback=slategray!8, colframe=slategray!50, arc=5pt, boxrule=1pt,
  left=12pt, right=12pt, top=8pt, bottom=8pt,
  title={\small\bfseries\color{white} #1},
  colbacktitle=slategray!70, coltitle=white,
  attach boxed title to top left={xshift=10pt, yshift=-\tcboxedtitleheight/2},
  boxed title style={arc=3pt, colframe=slategray!70, boxrule=0pt},
}

\geometry{margin=0.9in, top=1in, bottom=1in}

% ── Colour palette (Deep Ocean) ─────────────────────────
\definecolor{primary}  {RGB}{10,  80,  90}   % deep teal
\definecolor{accent}   {RGB}{0,  175, 185}   % cyan
\definecolor{lightbg}  {RGB}{224,247,250}    % pale cyan
\definecolor{jade}     {RGB}{0,  130,  80}
\definecolor{jadebg}   {RGB}{215,245,230}
\definecolor{amberclr} {RGB}{160, 90,   0}
\definecolor{amberbg}  {RGB}{255,243,220}
\definecolor{purpleclr}{RGB}{80,  30, 130}
\definecolor{purplebg} {RGB}{238,228,255}
\definecolor{redclr}   {RGB}{160,  30,  30}
\definecolor{redbg}    {RGB}{255,228,228}
\definecolor{greenclr} {RGB}{20, 110,  50}
\definecolor{greenbg}  {RGB}{215,248,228}
\definecolor{grayclr}  {RGB}{75,  85,  95}
\definecolor{graybg}   {RGB}{244,246,248}
\definecolor{darktext} {RGB}{25,  40,  50}
\definecolor{slategray}{RGB}{90, 105, 115}
\definecolor{rust}     {RGB}{180, 70,  20}
\definecolor{rustbg}   {RGB}{255,235,220}

% ── TOC styling ─────────────────────────────────────────
\setcounter{secnumdepth}{0}
\setcounter{tocdepth}{1}
\renewcommand{\contentsname}{}
\titlecontents{section}
  [18pt]
  {\addvspace{5pt}\bfseries\color{primary}}
  {}{}
  {\color{accent!70}\titlerule*[5pt]{.}\textbf{\color{accent}\contentspage}}

% ── Page style ──────────────────────────────────────────
\pagestyle{fancy}
\fancyhf{}
\fancyhead[L]{\small\color{primary}\textbf{Fluid Dynamics I}\;\textcolor{accent}{|}\;Complete Notes}
\fancyhead[R]{\small\color{darktext}BS-IV Mathematics}
\renewcommand{\headrulewidth}{0pt}
\fancyhead[C]{\color{accent}\rule[-1ex]{\textwidth}{1.2pt}}
\fancyfoot[C]{\small\color{darktext}\thepage}

% ── Section styling ─────────────────────────────────────
\titleformat{\section}
  {\large\bfseries\color{primary}}{}
  {0em}{}
  [\color{accent}\rule{\linewidth}{2pt}]
\titlespacing*{\section}{0pt}{20pt}{10pt}

\titleformat{\subsection}
  {\normalsize\bfseries\color{accent}}{}
  {0em}{}
\titlespacing*{\subsection}{0pt}{12pt}{4pt}

% ── tcolorbox environments ──────────────────────────────
\newtcolorbox{defbox}[1][]{
  enhanced, arc=4pt, boxrule=1.5pt,
  colback=jadebg, colframe=jade,
  left=10pt, right=10pt, top=8pt, bottom=8pt,
  title={\small\bfseries\color{white} Definition\ifx&#1&\else\ ---\ #1\fi},
  colbacktitle=jade, coltitle=white,
  attach boxed title to top left={xshift=10pt, yshift=-\tcboxedtitleheight/2},
  boxed title style={arc=3pt, colframe=jade, boxrule=0pt},
}

\newtcolorbox{thmbox}[1][]{
  enhanced, arc=4pt, boxrule=1.5pt,
  colback=amberbg, colframe=amberclr,
  left=10pt, right=10pt, top=8pt, bottom=8pt,
  title={\small\bfseries\color{white} Theorem\ifx&#1&\else\ ---\ #1\fi},
  colbacktitle=amberclr, coltitle=white,
  attach boxed title to top left={xshift=10pt, yshift=-\tcboxedtitleheight/2},
  boxed title style={arc=3pt, colframe=amberclr, boxrule=0pt},
}

\newtcolorbox{derivbox}[1][]{
  enhanced, arc=4pt, boxrule=1.5pt,
  colback=lightbg, colframe=primary,
  left=10pt, right=10pt, top=8pt, bottom=8pt,
  title={\small\bfseries\color{white} Derivation\ifx&#1&\else\ ---\ #1\fi},
  colbacktitle=primary, coltitle=white,
  attach boxed title to top left={xshift=10pt, yshift=-\tcboxedtitleheight/2},
  boxed title style={arc=3pt, colframe=primary, boxrule=0pt},
}

\newtcolorbox{keybox}[1][]{
  enhanced,
  borderline west={4pt}{0pt}{purpleclr},
  colback=purplebg!60, colframe=white,
  arc=0pt, boxrule=0pt,
  left=12pt, right=8pt, top=6pt, bottom=6pt,
}

\newtcolorbox{formulabox}{
  enhanced, arc=4pt, boxrule=1.5pt,
  colback=lightbg, colframe=accent,
  left=10pt, right=10pt, top=8pt, bottom=8pt,
  title={\small\bfseries\color{white} Key Formula},
  colbacktitle=accent, coltitle=white,
  attach boxed title to top left={xshift=10pt, yshift=-\tcboxedtitleheight/2},
  boxed title style={arc=3pt, colframe=accent, boxrule=0pt},
}

\newtcolorbox{notebox}{
  enhanced,
  borderline west={4pt}{0pt}{redclr},
  colback=redbg!60, colframe=white,
  arc=0pt, boxrule=0pt,
  left=12pt, right=8pt, top=6pt, bottom=6pt,
}

\newcolumntype{C}[1]{>{\centering\arraybackslash}p{#1}}
\newcommand{\hcell}[1]{\textbf{\color{white}#1}}

\begin{document}

% ════════════════════════════════════════════════════════════
%  COVER PAGE
% ════════════════════════════════════════════════════════════
\thispagestyle{empty}

\begin{tcolorbox}[
  enhanced, arc=0pt, boxrule=0pt,
  colback=primary, colframe=primary,
  left=20pt, right=20pt, top=28pt, bottom=20pt,
  borderline south={5pt}{0pt}{accent},
]
  \centering
  {\fontsize{26}{32}\bfseries\color{white}
    Fluid Dynamics I\\[8pt]
    Complete Notes
  }\\[12pt]
  {\large\color{lightbg}
    Fluid Statics \textbf{·} Kinematics \textbf{·} Dynamics \textbf{·} Bernoulli \textbf{·} Viscosity
  }
\end{tcolorbox}

\vspace{8pt}

\begin{tcolorbox}[
  enhanced, arc=0pt, boxrule=0pt,
  colback=white, colframe=white,
  borderline north={2pt}{0pt}{primary},
  borderline south={2pt}{0pt}{primary},
  left=0pt, right=0pt, top=0pt, bottom=0pt,
]
\begin{tabular}{p{0.33\textwidth} | p{0.33\textwidth} | p{0.31\textwidth}}
  \hspace{12pt}\begin{minipage}[t]{0.3\textwidth}\vspace{8pt}
    {\footnotesize\bfseries\color{accent} DEPARTMENT}\\[4pt]
    {\bfseries\color{primary} Mathematics}\\[3pt]
    {\footnotesize\color{darktext} SALU, Khairpur}
  \vspace{8pt}\end{minipage}
  &
  \hspace{12pt}\begin{minipage}[t]{0.3\textwidth}\vspace{8pt}
    {\footnotesize\bfseries\color{accent} COMPOSED BY}\\[4pt]
    {\bfseries\color{primary} Nadir Ali Khan}\\[3pt]
    {\footnotesize\color{darktext} BS-IV Mathematics}
  \vspace{8pt}\end{minipage}
  &
  \hspace{12pt}\begin{minipage}[t]{0.28\textwidth}\vspace{8pt}
    {\footnotesize\bfseries\color{accent} SESSION}\\[4pt]
    {\bfseries\color{primary} 2026}\\[3pt]
    {\footnotesize\color{darktext} Seat No: MT-0123-052}
  \vspace{8pt}\end{minipage}
\end{tabular}
\end{tcolorbox}

\vspace{16pt}

\begin{tcolorbox}[
  enhanced, arc=4pt, boxrule=1.5pt,
  colback=white, colframe=primary,
  left=10pt, right=10pt, top=4pt, bottom=8pt,
  title={\bfseries\color{white} Table of Contents},
  colbacktitle=primary,
  attach boxed title to top left={xshift=10pt, yshift=-\tcboxedtitleheight/2},
  boxed title style={arc=3pt, colframe=primary, boxrule=0pt},
]
\vspace{4pt}
\tableofcontents
\end{tcolorbox}

\newpage

% ════════════════════════════════════════════════════════════
\section{1. Introduction to Fluid Dynamics}
% ════════════════════════════════════════════════════════════

\begin{defbox}[Fluid]
A \textbf{fluid} is a substance that continuously deforms under an applied shear stress, no matter how small. Liquids and gases are both fluids.
\end{defbox}

\subsection{Types of Fluids}
\begin{itemize}[leftmargin=16pt, itemsep=3pt]
  \item \textbf{Ideal (Perfect) Fluid:} Incompressible, inviscid (no viscosity)
  \item \textbf{Real Fluid:} Has viscosity and compressibility
  \item \textbf{Newtonian Fluid:} Shear stress $\propto$ rate of strain (e.g.\ water, air)
  \item \textbf{Non-Newtonian Fluid:} Does not follow Newton's law of viscosity
  \item \textbf{Compressible:} Density changes with pressure (gases)
  \item \textbf{Incompressible:} Density constant (liquids)
\end{itemize}

\subsection{Key Fluid Properties}

\begin{center}
\renewcommand{\arraystretch}{1.4}
\begin{tabular}{|C{4cm}|C{5cm}|C{4cm}|}
\hline
\rowcolor{primary}
\hcell{Property} & \hcell{Formula} & \hcell{Unit (SI)} \\
\hline
Density & $\rho = m/V$ & kg/m$^3$ \\
\rowcolor{lightbg}
Specific Weight & $\gamma = \rho g$ & N/m$^3$ \\
Specific Gravity & $S = \rho_f / \rho_w$ & Dimensionless \\
\rowcolor{lightbg}
Dynamic Viscosity & $\mu = \tau / (du/dy)$ & Pa$\cdot$s \\
Kinematic Viscosity & $\nu = \mu/\rho$ & m$^2$/s \\
\rowcolor{lightbg}
Pressure & $P = F/A$ & Pa (N/m$^2$) \\
\hline
\end{tabular}
\end{center}

% ════════════════════════════════════════════════════════════
\section{2. Fluid Statics}
% ════════════════════════════════════════════════════════════

\begin{defbox}[Fluid Statics]
Study of fluids at \textbf{rest}. No relative motion between fluid layers, so no shear stress acts.
\end{defbox}

\subsection{Hydrostatic Pressure}

\begin{formulabox}
\[
  P = P_0 + \rho g h
\]
where $P_0$ = pressure at surface, $h$ = depth below surface.
\end{formulabox}

\subsection{Pascal's Law}

\begin{thmbox}[Pascal's Law]
Pressure applied to a confined fluid is transmitted equally in all directions throughout the fluid.
\[
  \frac{F_1}{A_1} = \frac{F_2}{A_2}
\]
\end{thmbox}

\subsection{Archimedes' Principle}

\begin{thmbox}[Archimedes' Principle]
A body wholly or partially submerged in a fluid experiences an upward buoyant force equal to the weight of fluid displaced:
\[
  F_B = \rho_f \, g \, V_{\text{displaced}}
\]
\end{thmbox}

\begin{itemize}[leftmargin=16pt]
  \item $F_B > W$ $\Rightarrow$ body floats
  \item $F_B = W$ $\Rightarrow$ body is neutrally buoyant
  \item $F_B < W$ $\Rightarrow$ body sinks
\end{itemize}

% ════════════════════════════════════════════════════════════
\section{3. Kinematics of Fluid Motion}
% ════════════════════════════════════════════════════════════

\begin{defbox}[Kinematics]
Study of fluid motion without considering the forces causing it.
\end{defbox}

\subsection{Types of Flow}

\begin{itemize}[leftmargin=16pt, itemsep=4pt]
  \item \textbf{Steady Flow:} Fluid properties at any point do not change with time $\left(\dfrac{\partial}{\partial t} = 0\right)$
  \item \textbf{Unsteady Flow:} Properties change with time
  \item \textbf{Uniform Flow:} Velocity same at every point at a given instant
  \item \textbf{Non-uniform Flow:} Velocity varies from point to point
  \item \textbf{Laminar Flow:} Fluid layers slide smoothly (low Reynolds number, Re $< 2000$)
  \item \textbf{Turbulent Flow:} Chaotic mixing of layers (Re $> 4000$)
  \item \textbf{Rotational Flow:} Fluid particles rotate about their own axis
  \item \textbf{Irrotational Flow:} No rotation; curl of velocity = 0
\end{itemize}

\subsection{Reynolds Number}

\begin{formulabox}
\[
  \text{Re} = \frac{\rho V D}{\mu} = \frac{VD}{\nu}
\]
Re $< 2000$: Laminar \quad Re $> 4000$: Turbulent \quad $2000 <$ Re $< 4000$: Transitional
\end{formulabox}

\subsection{Velocity and Acceleration}

For a velocity field $\mathbf{V} = u\,\hat{i} + v\,\hat{j} + w\,\hat{k}$:

\begin{formulabox}
\textbf{Material (Substantial) Derivative:}
\[
  \frac{D\mathbf{V}}{Dt} = \frac{\partial \mathbf{V}}{\partial t} + (\mathbf{V}\cdot\nabla)\mathbf{V}
\]
\[
  a_x = \frac{Du}{Dt} = \frac{\partial u}{\partial t} + u\frac{\partial u}{\partial x} + v\frac{\partial u}{\partial y} + w\frac{\partial u}{\partial z}
\]
\end{formulabox}

\subsection{Streamlines, Pathlines, Streaklines}

\begin{itemize}[leftmargin=16pt, itemsep=3pt]
  \item \textbf{Streamline:} Curve tangent to velocity vector at every point at a given instant
  \item \textbf{Pathline:} Actual path traced by a fluid particle over time
  \item \textbf{Streakline:} Locus of particles that have passed through a given point
\end{itemize}

\begin{keybox}
In \textbf{steady flow}, streamlines = pathlines = streaklines.
\end{keybox}

\subsection{Equation of a Streamline}
\begin{formulabox}
\[
  \frac{dx}{u} = \frac{dy}{v} = \frac{dz}{w}
\]
\end{formulabox}

% ════════════════════════════════════════════════════════════
\section{4. Equation of Continuity}
% ════════════════════════════════════════════════════════════

\begin{defbox}[Continuity Equation]
Based on conservation of mass: mass can neither be created nor destroyed.
\end{defbox}

\subsection{General Form (3D, Compressible)}

\begin{formulabox}
\[
  \frac{\partial \rho}{\partial t} + \frac{\partial(\rho u)}{\partial x} + \frac{\partial(\rho v)}{\partial y} + \frac{\partial(\rho w)}{\partial z} = 0
\]
Or in vector form: $\quad \dfrac{\partial \rho}{\partial t} + \nabla \cdot (\rho \mathbf{V}) = 0$
\end{formulabox}

\subsection{For Incompressible Flow ($\rho = \text{const}$)}

\begin{formulabox}
\[
  \frac{\partial u}{\partial x} + \frac{\partial v}{\partial y} + \frac{\partial w}{\partial z} = 0
  \qquad \text{i.e.} \quad \nabla \cdot \mathbf{V} = 0
\]
\end{formulabox}

\subsection{Continuity for 1D Pipe Flow}

\begin{formulabox}
\[
  A_1 V_1 = A_2 V_2 = Q \quad \text{(volume flow rate, m}^3\text{/s)}
\]
\end{formulabox}

\begin{examplebox*}[Worked Example]
Water flows through a pipe of diameter 0.2\,m at 3\,m/s. Find velocity where diameter narrows to 0.1\,m.
\[
  A_1V_1 = A_2V_2 \implies \frac{\pi(0.2)^2}{4}(3) = \frac{\pi(0.1)^2}{4}V_2 \implies V_2 = 12 \text{ m/s}
\]
\end{examplebox*}

% ════════════════════════════════════════════════════════════
\section{5. Euler's and Bernoulli's Equations}
% ════════════════════════════════════════════════════════════

\subsection{Euler's Equation of Motion}

\begin{derivbox}[Euler's Equation]
For an ideal (inviscid) fluid along a streamline:
\[
  \frac{\partial V}{\partial t} + V\frac{\partial V}{\partial s} = -\frac{1}{\rho}\frac{\partial P}{\partial s} - g\frac{\partial z}{\partial s}
\]
In vector form:
\[
  \rho\frac{D\mathbf{V}}{Dt} = -\nabla P + \rho\,\mathbf{g}
\]
\end{derivbox}

\subsection{Bernoulli's Equation}

\begin{thmbox}[Bernoulli's Equation]
For steady, incompressible, inviscid flow along a streamline:
\[
  \boxed{P + \frac{1}{2}\rho V^2 + \rho g z = \text{constant}}
\]
Or between two points:
\[
  P_1 + \frac{1}{2}\rho V_1^2 + \rho g z_1 = P_2 + \frac{1}{2}\rho V_2^2 + \rho g z_2
\]
\end{thmbox}

\subsection{Bernoulli Terms}

\begin{center}
\renewcommand{\arraystretch}{1.4}
\begin{tabular}{|C{3.5cm}|C{4cm}|C{5cm}|}
\hline
\rowcolor{primary}
\hcell{Term} & \hcell{Expression} & \hcell{Meaning} \\
\hline
Static Pressure & $P$ & Actual fluid pressure \\
\rowcolor{lightbg}
Dynamic Pressure & $\frac{1}{2}\rho V^2$ & Pressure due to motion \\
Hydrostatic Pressure & $\rho g z$ & Pressure due to elevation \\
\rowcolor{lightbg}
Total (Stagnation) & $P + \frac{1}{2}\rho V^2 + \rho gz$ & Constant along streamline \\
\hline
\end{tabular}
\end{center}

\begin{notebox}
\textbf{Assumptions for Bernoulli:} Steady flow $\cdot$ Incompressible $\cdot$ Inviscid $\cdot$ Along a streamline $\cdot$ No energy added or removed
\end{notebox}

\begin{examplebox*}[Worked Example — Bernoulli]
Water flows from a tank through a pipe. At point 1: $V_1 = 2$ m/s, $z_1 = 5$ m, $P_1 = 100$ kPa. At point 2: $z_2 = 0$, $V_2 = 4$ m/s. Find $P_2$.
\[
  P_2 = P_1 + \rho g(z_1-z_2) + \tfrac{1}{2}\rho(V_1^2 - V_2^2)
      = 100000 + 1000(9.81)(5) + \tfrac{1}{2}(1000)(4-16)
      = 143050 \text{ Pa}
\]
\end{examplebox*}

% ════════════════════════════════════════════════════════════
\section{6. Applications of Bernoulli's Equation}
% ════════════════════════════════════════════════════════════

\subsection{Venturi Meter}

Used to measure flow rate in a pipe.

\begin{formulabox}
\[
  Q = A_2\sqrt{\frac{2(P_1 - P_2)}{\rho\!\left(1 - (A_2/A_1)^2\right)}}
\]
\end{formulabox}

\subsection{Torricelli's Theorem}

Speed of fluid exiting a tank through a small hole at depth $h$:

\begin{formulabox}
\[
  V = \sqrt{2gh}
\]
\end{formulabox}

\subsection{Pitot Tube}

Measures fluid velocity using stagnation pressure:

\begin{formulabox}
\[
  V = \sqrt{\frac{2(P_0 - P)}{\rho}}
\]
where $P_0$ = stagnation pressure, $P$ = static pressure.
\end{formulabox}

% ════════════════════════════════════════════════════════════
\section{7. Viscosity and Viscous Flow}
% ════════════════════════════════════════════════════════════

\begin{defbox}[Viscosity]
\textbf{Viscosity} is the measure of a fluid's resistance to deformation under shear stress.
\end{defbox}

\subsection{Newton's Law of Viscosity}

\begin{formulabox}
\[
  \tau = \mu \frac{du}{dy}
\]
where $\tau$ = shear stress (Pa), $\mu$ = dynamic viscosity (Pa$\cdot$s), $\dfrac{du}{dy}$ = velocity gradient.
\end{formulabox}

\subsection{Navier-Stokes Equations}

For incompressible viscous flow:

\begin{formulabox}
\[
  \rho\frac{D\mathbf{V}}{Dt} = -\nabla P + \mu\,\nabla^2\mathbf{V} + \rho\,\mathbf{g}
\]
$x$-component:
\[
  \rho\!\left(\frac{\partial u}{\partial t} + u\frac{\partial u}{\partial x} + v\frac{\partial u}{\partial y}\right)
  = -\frac{\partial P}{\partial x} + \mu\!\left(\frac{\partial^2 u}{\partial x^2}+\frac{\partial^2 u}{\partial y^2}\right)
\]
\end{formulabox}

\subsection{Hagen-Poiseuille Flow (Pipe Flow)}

Steady laminar flow in a circular pipe of radius $R$:

\begin{formulabox}
\[
  u(r) = \frac{1}{4\mu}\!\left(-\frac{dP}{dx}\right)(R^2 - r^2)
\]
\[
  Q = \frac{\pi R^4}{8\mu}\!\left(-\frac{dP}{dx}\right) \qquad
  V_{\max} = \frac{R^2}{4\mu}\!\left(-\frac{dP}{dx}\right)
\]
\end{formulabox}

\begin{keybox}
In Hagen-Poiseuille flow: $V_{\text{avg}} = \dfrac{V_{\max}}{2}$. The velocity profile is parabolic.
\end{keybox}

% ════════════════════════════════════════════════════════════
\section{8. Irrotational Flow and Stream Function}
% ════════════════════════════════════════════════════════════

\subsection{Vorticity and Irrotational Flow}

\begin{formulabox}
\textbf{Vorticity:} $\boldsymbol{\omega} = \nabla \times \mathbf{V}$

For 2D flow: $\omega_z = \dfrac{\partial v}{\partial x} - \dfrac{\partial u}{\partial y}$

\textbf{Irrotational:} $\omega_z = 0 \implies \dfrac{\partial v}{\partial x} = \dfrac{\partial u}{\partial y}$
\end{formulabox}

\subsection{Stream Function $\psi$}

\begin{defbox}[Stream Function]
A scalar function $\psi(x,y)$ such that:
\[
  u = \frac{\partial \psi}{\partial y}, \qquad v = -\frac{\partial \psi}{\partial x}
\]
Automatically satisfies continuity. Streamlines are curves $\psi = \text{const}$.
\end{defbox}

\begin{formulabox}
Flow rate between two streamlines:
\[
  Q = \psi_2 - \psi_1
\]
\end{formulabox}

\subsection{Velocity Potential $\phi$}

\begin{defbox}[Velocity Potential]
For irrotational flow, $\exists$ scalar $\phi$ such that:
\[
  u = \frac{\partial \phi}{\partial x}, \qquad v = \frac{\partial \phi}{\partial y}, \qquad \mathbf{V} = \nabla\phi
\]
Equipotential lines are curves $\phi = \text{const}$ and are \textbf{orthogonal} to streamlines.
\end{defbox}

\subsection{Laplace's Equation}

For irrotational, incompressible flow:
\begin{formulabox}
\[
  \nabla^2\phi = \frac{\partial^2\phi}{\partial x^2} + \frac{\partial^2\phi}{\partial y^2} = 0
  \qquad \text{and} \qquad
  \nabla^2\psi = 0
\]
\end{formulabox}

% ════════════════════════════════════════════════════════════
\section{9. Dimensional Analysis}
% ════════════════════════════════════════════════════════════

\subsection{Buckingham $\pi$ Theorem}

\begin{thmbox}[Buckingham $\pi$ Theorem]
If a physical problem involves $n$ variables and $m$ fundamental dimensions (M, L, T), then it can be expressed in terms of $(n - m)$ dimensionless $\pi$ groups.
\end{thmbox}

\subsection{Important Dimensionless Numbers}

\begin{center}
\renewcommand{\arraystretch}{1.5}
\begin{tabular}{|C{3.5cm}|C{4cm}|C{5cm}|}
\hline
\rowcolor{primary}
\hcell{Number} & \hcell{Formula} & \hcell{Physical Meaning} \\
\hline
Reynolds (Re) & $\rho VL/\mu$ & Inertia / Viscous forces \\
\rowcolor{lightbg}
Froude (Fr) & $V/\sqrt{gL}$ & Inertia / Gravity forces \\
Euler (Eu) & $P/\rho V^2$ & Pressure / Inertia forces \\
\rowcolor{lightbg}
Mach (Ma) & $V/c$ & Flow / Sound speed \\
Weber (We) & $\rho V^2 L/\sigma$ & Inertia / Surface tension \\
\hline
\end{tabular}
\end{center}

% ════════════════════════════════════════════════════════════
\section{Quick Reference Cheat Sheet}
% ════════════════════════════════════════════════════════════

\begin{tcolorbox}[
  enhanced, arc=0pt, boxrule=0pt,
  colback=primary, colframe=primary,
  left=14pt, right=14pt, top=10pt, bottom=10pt,
  borderline south={3pt}{0pt}{accent},
]
  \centering
  {\large\bfseries\color{white} FLUID DYNAMICS I — MASTER FORMULA SHEET}\\[4pt]
  {\small\color{lightbg} All key equations and results for the exam}
\end{tcolorbox}

\vspace{10pt}

\begin{tcolorbox}[
  enhanced, arc=4pt, boxrule=1.5pt,
  colback=jadebg, colframe=jade,
  left=10pt, right=10pt, top=8pt, bottom=8pt,
  title={\small\bfseries\color{white} Core Equations},
  colbacktitle=jade, coltitle=white,
  attach boxed title to top left={xshift=10pt, yshift=-\tcboxedtitleheight/2},
  boxed title style={arc=3pt, colframe=jade, boxrule=0pt},
  width=0.48\linewidth, nobeforeafter,
]
\renewcommand{\arraystretch}{1.5}
\begin{tabular}{@{}ll@{}}
  Hydrostatic P & $P = P_0 + \rho gh$ \\
  Buoyancy & $F_B = \rho_f g V$ \\
  Continuity & $A_1V_1 = A_2V_2$ \\
  Bernoulli & $P+\tfrac{1}{2}\rho V^2+\rho gz=C$ \\
  Torricelli & $V = \sqrt{2gh}$ \\
  Reynolds & $\text{Re} = \rho VD/\mu$ \\
  Viscosity & $\tau = \mu\,du/dy$ \\
  H-P Flow & $Q = \pi R^4 (-dP/dx)/8\mu$ \\
\end{tabular}
\end{tcolorbox}
\hfill
\begin{tcolorbox}[
  enhanced, arc=4pt, boxrule=1.5pt,
  colback=lightbg, colframe=primary,
  left=10pt, right=10pt, top=8pt, bottom=8pt,
  title={\small\bfseries\color{white} Stream Function \& Potential},
  colbacktitle=primary, coltitle=white,
  attach boxed title to top left={xshift=10pt, yshift=-\tcboxedtitleheight/2},
  boxed title style={arc=3pt, colframe=primary, boxrule=0pt},
  width=0.48\linewidth, nobeforeafter,
]
\renewcommand{\arraystretch}{1.5}
\begin{tabular}{@{}ll@{}}
  Stream fn & $u=\partial\psi/\partial y,\;v=-\partial\psi/\partial x$ \\
  Vel. potential & $u=\partial\phi/\partial x,\;v=\partial\phi/\partial y$ \\
  Laplace & $\nabla^2\phi = 0,\;\nabla^2\psi = 0$ \\
  Irrotational & $\partial v/\partial x = \partial u/\partial y$ \\
  Vorticity & $\omega_z = \partial v/\partial x - \partial u/\partial y$ \\
  Mat.\ Deriv. & $D/Dt = \partial/\partial t + \mathbf{V}\cdot\nabla$ \\
  Continuity & $\nabla\cdot\mathbf{V} = 0$ \\
  Euler & $\rho D\mathbf{V}/Dt = -\nabla P + \rho\mathbf{g}$ \\
\end{tabular}
\end{tcolorbox}

\vspace{10pt}

\begin{tcolorbox}[
  enhanced, arc=4pt, boxrule=1.5pt,
  colback=amberbg, colframe=amberclr,
  left=10pt, right=10pt, top=8pt, bottom=8pt,
  title={\small\bfseries\color{white} Dimensionless Numbers},
  colbacktitle=amberclr, coltitle=white,
  attach boxed title to top left={xshift=10pt, yshift=-\tcboxedtitleheight/2},
  boxed title style={arc=3pt, colframe=amberclr, boxrule=0pt},
]
\[
  \text{Re} = \frac{\rho VL}{\mu} \qquad
  \text{Fr} = \frac{V}{\sqrt{gL}} \qquad
  \text{Ma} = \frac{V}{c} \qquad
  \text{We} = \frac{\rho V^2 L}{\sigma} \qquad
  \text{Eu} = \frac{P}{\rho V^2}
\]
\[
  \text{Re} < 2000: \text{ Laminar} \qquad
  \text{Re} > 4000: \text{ Turbulent} \qquad
  2000 < \text{Re} < 4000: \text{ Transitional}
\]
\end{tcolorbox}

\vspace{6pt}

\begin{tcolorbox}[
  enhanced, arc=4pt, boxrule=1.5pt,
  colback=purplebg!60, colframe=purpleclr,
  left=10pt, right=10pt, top=8pt, bottom=8pt,
  title={\small\bfseries\color{white} Flow Types — Quick Recall},
  colbacktitle=purpleclr, coltitle=white,
  attach boxed title to top left={xshift=10pt, yshift=-\tcboxedtitleheight/2},
  boxed title style={arc=3pt, colframe=purpleclr, boxrule=0pt},
  width=0.48\linewidth, nobeforeafter,
]
\begin{itemize}[leftmargin=12pt, itemsep=2pt]
  \item Steady: $\partial/\partial t = 0$
  \item Uniform: $\partial/\partial s = 0$
  \item Incompressible: $\nabla\cdot\mathbf{V}=0$
  \item Irrotational: $\nabla\times\mathbf{V}=0$
  \item Laminar: Re $< 2000$
  \item Turbulent: Re $> 4000$
\end{itemize}
\end{tcolorbox}
\hfill
\begin{tcolorbox}[
  enhanced, arc=4pt, boxrule=1.5pt,
  colback=greenbg, colframe=greenclr,
  left=10pt, right=10pt, top=8pt, bottom=8pt,
  title={\small\bfseries\color{white} Bernoulli Applications},
  colbacktitle=greenclr, coltitle=white,
  attach boxed title to top left={xshift=10pt, yshift=-\tcboxedtitleheight/2},
  boxed title style={arc=3pt, colframe=greenclr, boxrule=0pt},
  width=0.48\linewidth, nobeforeafter,
]
\begin{itemize}[leftmargin=12pt, itemsep=2pt]
  \item Venturi meter: flow rate in pipe
  \item Pitot tube: fluid velocity
  \item Torricelli: tank exit velocity
  \item Orifice: discharge coefficient
  \item Aerofoil: lift generation
  \item Syphon: fluid transfer
\end{itemize}
\end{tcolorbox}

\vspace{20pt}
\begin{center}
  \color{grayclr}\rule{0.6\linewidth}{0.4pt}\\[8pt]
  {\small\color{darktext} For feedback or to request additional topics:}\\[4pt]
  {\normalsize\bfseries\color{primary} +92 304 2019543}
\end{center}

\end{document}
